%%% FIELD proof-prop-no-terminal-effect-reduction
Assume the antecedent \(D(0)=D(1)\) almost surely in the compatible full-law class.  This submodel satisfies the maintained monotonicity restriction \(D(0)\leq D(1)\).  Fix \(x\in\mathcal X\), \(t\in[0,\tau]\), and \(\nu_x\in\mathcal N_x\).  For every Borel set \(B\subseteq\mathcal E\), define the finite measures
\[
q_{x,t}(B)=\nu_x\bigl((t,\infty]\times B\bigr),
\qquad
k_{x,t}(B)=K_x(t,B)
=\mu_{1x}\{D(1)>t,\ R(1)\in B\}.
\]
By the survivor-capacity restriction in \(\mathcal N_x\), applied at the deadline \(u=t\),
\[
q_{x,t}(B)\leq k_{x,t}(B)
\quad\text{for every Borel }B\subseteq\mathcal E. \tag{1}
\]
The first marginal of \(\nu_x\) is \(d_x\), so
\[
q_{x,t}(\mathcal E)
=\nu_x\bigl((t,\infty]\times\mathcal E\bigr)
=d_x((t,\infty])
=\Pr\{D(0)>t\mid X=x\}. \tag{2}
\]
The definition of \(K_x\) gives
\[
k_{x,t}(\mathcal E)
=K_x(t,\mathcal E)
=\Pr\{D(1)>t\mid X=x\}. \tag{3}
\]
Because \(D(0)=D(1)\) almost surely, their conditional survivor probabilities agree in every positive-probability covariate cell.  Equations (2)--(3) therefore imply
\[
q_{x,t}(\mathcal E)=k_{x,t}(\mathcal E). \tag{4}
\]
The finite-measure domination in (1) makes \(k_{x,t}-q_{x,t}\) a nonnegative finite measure.  Its total mass is zero by (4); hence, for every Borel \(B\),
\[
0\leq (k_{x,t}-q_{x,t})(B)
\leq (k_{x,t}-q_{x,t})(\mathcal E)=0.
\]
Thus \(q_{x,t}=k_{x,t}\).  Taking \(B=[0,r]\) proves, for every \(0\leq r\leq t\leq\tau\),
\[
\nu_x\{D(0)>t,\ R(1)\leq r\}
=\mu_{1x}\{D(1)>t,\ R(1)\leq r\}. \tag{5}
\]

Now let \((r,t)\in\mathcal H_\eta\).  The target-domain definition gives \(S(t)\geq\eta>0\), so division by \(S(t)\) is valid.  Substituting (5) and the definition
\(b_x(r,t)=\mu_{0x}\{D(0)>t,R(0)\leq r\}\) into the TV-SACE functional yields
\[
\theta_\nu(r,t)
=\frac{1}{S(t)}\sum_{x\in\mathcal X}p_x
\left[
\mu_{1x}\{D(1)>t,R(1)\leq r\}
-\mu_{0x}\{D(0)>t,R(0)\leq r\}
\right]
=:\bar\theta(r,t). \tag{6}
\]
The right-hand side contains no coupling \(\nu_x\).  Moreover, the within-arm laws in (6) are functionals of \(P\) by the IPCW-recovery result used in the capacity-projection theorem; \(p_x\) and \(S(t)\) are also functionals of \(P\).  Thus \(\bar\theta\) is an observed-law identifying functional on \(\mathcal H_\eta\), rather than a definition of the causal target.

The connection to the causal target follows from the no-terminal-effect premise.  Indeed,
\[
\{D(0)>t,D(1)>t\}=\{D(0)>t\}=\{D(1)>t\}
\quad\text{almost surely}, \tag{7}
\]
and standardization over the finite covariate support gives
\[
S(t)=\Pr\{D(0)>t,D(1)>t\}. \tag{8}
\]
Applying (7)--(8) to the two standardized numerators in (6) proves
\[
\bar\theta(r,t)
=\Pr\{R(1)\leq r\mid D(0)>t,D(1)>t\}
-\Pr\{R(0)\leq r\mid D(0)>t,D(1)>t\}. \tag{9}
\]
Hence (6) is precisely the ordinary standardized within-arm cumulative-incidence contrast in the common survivor cohort.

For \(0\leq a\leq t\leq\tau\) with \(S(t)\geq\eta\), every \((r,t)\) with \(0\leq r\leq a\) belongs to \(\mathcal H_\eta\).  The RM-SACE definition and (6) consequently give
\[
\rho_\nu(a,t)
=-\int_0^a\theta_\nu(r,t)\,dr
=-\int_0^a\bar\theta(r,t)\,dr
=:\bar\rho(a,t), \tag{10}
\]
which also contains no coupling.  To verify its causal interpretation rather than infer it from the notation, use the pointwise identity, with \(\min\{\infty,a\}=a\),
\[
\min\{u,a\}=a-\int_0^a\mathbf 1\{u\leq r\}\,dr,
\qquad u\in\mathcal E. \tag{11}
\]
Integrating the bounded nonnegative indicator in (11) over each conditional event-time law and subtracting the two arms shows from (9)--(10) that
\[
\bar\rho(a,t)
=E\!\left[\min\{R(1),a\}-\min\{R(0),a\}
\mid D(0)>t,D(1)>t\right]. \tag{12}
\]

Finally, the sharp capacity-projection theorem states that every compatible causal surface pair is represented by a family \(\nu_x\in\mathcal N_x\), and every such family has a compatible full-law lift.  Equations (6) and (10) show that every represented pair equals \((\bar\theta,\bar\rho)\).  The compatible full law posited in the antecedent supplies at least one represented family, so the image is nonempty.  Therefore
\[
\Theta_I(P)=\{(\bar\theta,\bar\rho)\},
\]
and (9) and (12) identify this singleton with the ordinary standardized TV-SACE and RM-SACE contrasts in the common survivor cohort.
